
\section{Four steps for presenting a dynamic-programming algorithm}\label{dp: sec: four steps}\doHeader{Four steps for presenting DP algorithms}

In the rest of these notes,
we use the following four steps to present each dynamic-programming algorithm:

\begin{description}
\item[1. Define the subproblems.]
  Define the subproblems to be solved for a given instance,
  and state how the final answer is computed from these subproblems.
  For example, to present the dynamic-programming algorithm for \prob{Making Change},
  the first step would be as follows:

  \medskip 

  \emph{1.   Given the instance $(D, T)$ of \prob{Making Change},
    define the following subproblems.
    For each $t\in \{0,1,\ldots,T\}$,
    define $M(t)$ to be the minimum number of coins needed to make change
    (using denominations in $D$) of total value $t$.
    The solution to the problem is $M(T)$.}

  \medskip 

  This tells the reader exactly what subproblems the algorithm will solve,
  and also how (once they are solved) the algorithm will compute the output for the given instance.

\item[2. State the recurrence relation.]  The recurrence relation is a divide-and-conquer equation
  that can be used to compute the value for each subproblem
  from the values of some smaller subproblems.
  The left-hand side of the equation is the value of the subproblem.
  The right-hand side is an expression
  built from the values of smaller subproblems.
  For \prob{Making Change}, the recurrence relation is in Lemma~\ref{dp: lemma: change}:

 \medskip 
  
    \emph{2. 
      \(
      M(t) = \begin{cases}
        0 & \text{ if } t=0 \\
        \min\{ 1 + M(t-d_i) : d_i \in D,\, d_i \le t\} & \text{ otherwise}.
      \end{cases}
      \)}

  \medskip 

  The recurrence relation does not \emph{define} $M(t)$!
  Rather, it is an identity that holds for $M(t)$ as $M$ is defined in the first step.
  (To prove correctness of the algorithm, we'll have to prove that the recurrence relation holds.)
  
  The first and second step together determine the algorithm:
  the algorithm iterate (or recurses) over the set of subproblems
  to be solved.  It solves each subproblem using the recurrence relation,
  evaluating the right-hand side to determine the value of the left-hand side.

  For readers who are familiar with dynamic programming,
  the first two steps are enough to implicitly define the algorithm,
  so it is usually unnecessary to give any further definition.
  
\item[3. Prove correctness.]
  The third step is to prove that (or explain why) the recurrence relation holds. 
  As long as it holds,
  correctness of the algorithm (iterative, or recursive and memoized)
  will follow by induction.
  
\item[4. Bound the run time.] Analyze the runtime of the algorithm.
  This is the sum, over all of the subproblems,
  of the time to evaluate
  the right-hand side of the recurrence relation
  (assuming the smaller subproblems are already solved).
  For \prob{Making Change}, this would be enough:
  
  \medskip

  \emph{
    4. There are $T+1$ subproblems (one for each $t\in\{0,1,\ldots,T\}$;
    subproblem $t$ is to compute $M(t)$).
    For each, the time to evaluate the right-hand side of the recurrence for $M(t)$ is $O(k)$
    (the number of denominations).
    So the total time for the algorithm is $O(T k)$.
  }
\end{description}

%%% Local Variables:
%%% mode: latex
%%% TeX-master: "dynamic_programming"
%%% End:

%%% Local Variables:
%%% mode: latex
%%% TeX-master: "dynamic_programming"
%%% End:
